\section{STALE as a State-Validity Case Study}
\label{sec:stale}

STALE narrows long-term memory to a specific memory-use subtask: determining whether a later observation implicitly invalidates an earlier memory \cite{chao2026stalellmagentsknow}.
This setting is valuable because the failure is not ordinary fact retrieval.
The system must revise state, resist stale premises, and adapt actions to the currently authorized state.

\subsection{Primitive-Role Alignment}
\label{subsec:stale-alignment}

In our terminology, stale-state validity requires a state-oriented artifact substrate and a primitive stack that can maintain and expose current authority.
A stylized primitive set is
\begin{align}
\primitiveSet_{\mathrm{STALE}}=\{&
\rho_{\mathrm{extract}}, \rho_{\mathrm{propagate}},
\rho_{\mathrm{adjudicate}}, \nonumber\\
&\rho_{\mathrm{authorize}}, \rho_{\mathrm{readout}}\}.
\end{align}
These primitives correspond to extracting state changes, propagating implications through a dependency structure, adjudicating conflicts, authorizing the current state, and reading out the evidence needed for a query.

The decision-facing path is:
\begin{equation}
\artifactstore_t \rightarrow \artifactstore_t^{\mathrm{auth}} \rightarrow v_u^{\mathrm{state}} \rightarrow \hat{y}.
\end{equation}
This path makes explicit why retrieval visibility is insufficient.
The model must know which state is current, why an older premise is stale, and how this should govern the answer.

\subsection{Probe-Level Interpretation}
\label{subsec:stale-probes}

STALE uses three probe families:
\begin{itemize}[leftmargin=*]
    \item \textbf{State Resolution:} whether the system identifies the current state.
    \item \textbf{Premise Resistance:} whether the system rejects or corrects a question that embeds a stale premise.
    \item \textbf{Implicit Policy Adaptation:} whether the system adapts recommendations or actions to the updated state.
\end{itemize}
These probes map naturally onto evidence-view roles: current authority, stale-premise marking, dependency explanation, and action-grounding support.

\subsection{Planned Reporting}
\label{subsec:stale-reporting}

The STALE section will report whether a specialized state artifact and primitive stack improves the exposure of current-validity roles, and which probe families benefit most.
If the experiments show limited gains, this section will be reframed as a limitation analysis rather than a positive method claim.

\begin{table}[t]
\centering
\small
\begin{tabular}{lccc}
\toprule
System/View & SR & PR & IPA \\
\midrule
Retrieval-only readout & -- & -- & -- \\
Authorized state readout & -- & -- & -- \\
Role-labeled state view & -- & -- & -- \\
\bottomrule
\end{tabular}
\caption{Planned STALE probe-level reporting. SR: State Resolution; PR: Premise Resistance; IPA: Implicit Policy Adaptation.}
\label{tab:stale-placeholder}
\end{table}

